\par\noindent\textbf{(1)} 总自旋 \(\boldsymbol S=\boldsymbol S_1+\boldsymbol S_2\)，有 \(\boldsymbol S_1\cdot\boldsymbol S_2=[S^2-S_1^2-S_2^2]/2\)。两单自旋均为 \(\hbar^2(3/4)\)，总自旋 \(S=0\) singlet给能 \(E_s=-3J/4\)，三重 \(S=1\) triplet给 \(E_t=J/4\)。
\par\noindent\textbf{(2)} 热配分函数 \(Z=e^{3\beta J/4}+3e^{-\beta J/4}\)，平均能量 \(U=[(-3J/4)e^{3\beta J/4}+3(J/4)e^{-\beta J/4}]/Z\)。能隙为 \(J\)，低温主要占据singlet，高温各四个自旋态趋等概率。
\par\noindent\textbf{(3)} 热态为singlet与三重子空间的各向同性混合，singlet概率 \(p_s=e^{\beta J}/(e^{\beta J}+3)\)。两qubit的Bell对角态当最大本征概率超过 \(1/2\) 时纠缠，因此 \(p_s>1/2\) 等价于 \(e^{\beta J}>3\)，即 \(k_BT<J/\ln3\)。
\par\noindent\textbf{(4)} 当 \(T\to0\)，\(p_s\to1\)，态趋纯singlet；当 \(T\to\infty\)，\(p_s\to1/4\)，密度矩阵趋 \(I_4/4\) 而不纠缠。这里的温度纠缠阈值是理想两自旋、各向同性交换模型结果，外场和各向异性会改变。
